\documentclass[12pt]{article}
\usepackage[utf8]{inputenc}
\usepackage[T1]{fontenc}
\usepackage{hyperref}
\title{Matsuva Language Resources for K-12 Curriculum Development}
\author{Compiled by Claude Sonnet 5 (claude-sonnet-5) for Andr\'as Kornai}
\date{}
\begin{document}
\maketitle

\section{Introduction}
Matsuva (ISO 639-3 tsh) is listed in the sources I found under the name Tsuvan. Alternate names there include Matsuvan, Motsuvan, Terki, Telaki, Teleki and Tchede. It is a Chadic (Afro-Asiatic) language of Cameroon's Far North Province, spoken around the Teleki massif. The search results give differing speaker estimates: about 2,300 (a 2000 SIL figure) and a population of 4,800 for the Terki/Teleki people. Very little exists. I found no school materials.

\section{Primary Educational Resources}
I found no primers, textbooks or school programmes for this language.

\section{Linguistic and Reference Materials}
\begin{itemize}
\item Tsuvan language, Wikipedia (overview, names, classification): \url{https://en.wikipedia.org/wiki/Tsuvan_language}
\item A rapid appraisal language survey of Tsuvan (Mayo-Tsanaga division, Far North Province), SIL Cameroon archive. The search result showed only the title, and it appears to assess whether a written form and literacy project is worthwhile: \url{https://www.silcam.org/resources/archives/9192}
\item UNESCO World Atlas of Languages entry, Tsuvan in Cameroon: \url{https://en.wal.unesco.org/countries/cameroon/languages/tsuvan}
\end{itemize}

\section{Religious and Community Resources}
\begin{itemize}
\item Joshua Project language page for tsh (population and religion data only): \url{https://joshuaproject.net/index.php/languages/tsh}
\end{itemize}
I found no community organisations.

\section{Digital and Audio Resources}
\begin{itemize}
\item Global Recordings Network page for tsh. I did not confirm what recordings it holds: \url{https://globalrecordings.net/en/language/tsh}
\item Wikitongues page for tsh. The search summary says it describes the language as unwritten: \url{https://wikitongues.org/languages/tsh/}
\end{itemize}

\section{Curriculum Development Context}
The language appears to be unwritten, and I found no standard orthography. No native-language instruction is documented. A curriculum would have to start with orthography development and documentation, together with the community. The sources disagree on speaker numbers and on vitality: one summary reports EGIDS 6a, while the task description says only older adults speak it. Both should be checked locally.

\end{document}